% Historical source (removed by author request): references/_OLD_Übungsaufgaben/Übungsblatt 5 - Differenzialrechnung/Blatt5.tex:272
\begin{ExerciseBlock}
\ExerciseHeading{13}{Optimierung: Abstand zum Ursprung}{differential_calculus_distance_optimization}
Gegeben sei eine umgekehrte Parabel.

\begin{center}
	\CachedTikzFigure{content/resources/chapters/differential_calculus/tikz/distance_optimization.tex}
\end{center}

\begin{enumerate}[label=\alph*)]
	
	\item Bestimmen Sie die Funktionsgleichung $f(x)$, der in der Abbildung in Grün dargestellten Funktion.
	
	\item Bestimmen Sie die Funktion $A(x)$, die den Abstand zwischen dem Ursprung und der Funktion am Punkt $(x,f(x))$ beschreibt.
	
	\item Für ein $x_0 \in (0,1)$ nimmt die Funktion $A(x)$ ein lokales Minimum an. Bestimmen Sie $x_0$. 
	
	\item Bestimmen Sie den kürzesten Abstand $A(x_0)$ zwischen dem Ursprung und der Funktion $f(x)$.
\end{enumerate}	

\end{ExerciseBlock}
\UseExerciseSolution{differential_calculus_distance_optimization}
